\subsection{变长轨迹的访问测度与梯度累积}
\label{sec:prefix-measure-reduction}

固定第 $k$ 轮 rollout 分布。令轨迹 $\tau$ 的有效 response 长度为
$H(\tau)$，包含生成的 EOS，排除 prompt 和 padding；在预算终止时保留实际生成的 token。
假定 $0<\mathbb E H<\infty$。本书明确采用归一化的 token 访问测度
\begin{equation}
 d_k(A)=\frac{\mathbb E_{\tau\sim q_k}
 [\sum_{t=1}^{H(\tau)}\mathbf1\{s_t\in A\}]}{\mathbb E H}.
\end{equation}
这里 prompt 也按训练数据分布采样。长轨迹贡献更多被访问的 prefix，
不是先将每条轨迹的总贡献归一化为一。固定 rollout 后，长度和测度均不对 inner 参数求导。

令每个有效位置的损失或 score surrogate 为 $\ell_\theta(s,a)$。
独立 rollout 的经验目标及梯度为
\begin{align}
 \mathsf T_k&=\sum_{i=1}^{B}H_i,\qquad
 \widehat L_k^{\mathrm{tok}}(\theta)
 =\frac1{\mathsf T_k}\sum_{i=1}^{B}\sum_{t=1}^{H_i}
 \ell_\theta(s_{i,t},a_{i,t}),\\
 \nabla\widehat L_k^{\mathrm{tok}}
 &=\frac1{\mathsf T_k}\sum_{i,t}\nabla\ell_\theta(s_{i,t},a_{i,t}).
\end{align}
总体目标为 $\mathbb E_{d_k, q_k}\ell_\theta$。有限 batch 的随机分母
使上述比率估计量一般不是总体目标的严格无偏估计；在独立轨迹和可积条件下，
分子、分母的样本均值分别收敛，其比值一致收敛到总体目标。
token 在同一轨迹内相关，因此不能把 $\mathsf T_k$ 个位置视为独立统计样本来计算置信区间。

\begin{proposition}[访问测度与条件动作分布]
上述 token 访问测度对应的总体 prefix--action 联合分布可分解为
$d_k(s) q_k(a\mid s)$。对每条轨迹额外乘以 $1/H(\tau)$ 的序列平均，
一般改变这一条件动作分布。
\end{proposition}
\begin{proof}
在时刻 $t$ 产生动作之前，事件 $\{t\le H\}$ 和 prefix $s_t$
已由历史决定。因此对任意可积函数 $h$，条件期望给出
\begin{align}
 \mathbb E\sum_{t\ge1}\mathbf1\{t\le H\}h(s_t,a_t)
 &=\mathbb E\sum_{t\ge1}\mathbf1\{t\le H\}
 \mathbb E_{a\sim q_k(\cdot\mid s_t)}h(s_t,a)\\
 &=\mathbb E H\;\mathbb E_{s\sim d_k,a\sim q_k}h(s,a).
\end{align}
除以 $\mathbb E H$ 得到所述分解。预算终止只改变历史可测的存活事件，
不改变该条件期望步骤。

对于序列权重 $1/H$，长度取决于当前及未来生成动作，不能把该权重
从动作条件期望中提出。举例：初始 prefix 下两个动作概率各为 $1/2$；
动作 $a$ 立即 EOS，长度为一；动作 $b$ 后确定生成一个 EOS，长度为二。
按轨迹长度归一化后，初始位置的联合质量分别为 $1/2$ 和 $1/4$，
故给定初始 prefix 的动作概率变成 $2/3$ 与 $1/3$，不同于原策略。
\end{proof}

相比之下，先在每条序列内平均、再在 batch 平均得到
\begin{equation}
 \widehat L_k^{\mathrm{seq}}=
 \frac1B\sum_i\frac1{H_i}\sum_t\ell_\theta(s_{i,t},a_{i,t}).
\end{equation}
同一 token 在两个目标中的权重之比为
$(1/(BH_i))/(1/\mathsf T_k)=\bar H/H_i$。
例如两条轨迹长度为十和一百，token 平均给两条轨迹的总质量分别为
$1/11$ 与 $10/11$；序列平均给它们各 $1/2$。
这两种目标都可以被定义，但不能在不修改访问测度、条件动作分布和梯度推导的情况下互换。
若每个 prefix 可独立表示、使用完整条件期望且支持集相同，
只改变正的 prefix 权重不改变逐 prefix 的精确 Gibbs 解；
共享参数投影、有限样本和上述依赖未来长度的权重没有这一等价保证。

不除以长度的 batch 序列和 $B^{-1}\sum_i\sum_t\ell$，
则恰等于 $\bar H\widehat L_k^{\mathrm{tok}}$。
在同一冻结 batch 内它是正的常数倍，最优点不变，
但普通梯度尺度、裁剪行为及跨 outer 的相对尺度不同；不能声称 Adam 和裁剪后轨迹也相同。

\begin{proposition}[micro-batch 划分不变性]
将本轮有效位置划分为任意不交 micro-batch $\mathcal B_m$。
若每次 backward 使用
$L_m=\mathsf T_k^{-1}\sum_{(i,t)\in\mathcal B_m}\ell_{i,t}$，
只在所有 backward 完成后裁剪并执行一次 optimizer step，
则累积梯度与全 batch 的梯度相同，忽略浮点求和次序误差。
\end{proposition}
\begin{proof}
梯度对有限和线性，且冻结分母 $\mathsf T_k$ 与 inner 参数无关，故
\begin{equation}
 \sum_m\nabla L_m
 =\frac1{\mathsf T_k}\sum_m\sum_{(i,t)\in\mathcal B_m}\nabla\ell_{i,t}
 =\nabla\widehat L_k^{\mathrm{tok}}.
\end{equation}
若把每个 $L_m$ 改为各自 token 均值后直接累加，
每个位置的权重变成 $1/|\mathcal B_m|$，不再是 $1/\mathsf T_k$。
只有再乘 $|\mathcal B_m|/\mathsf T_k$ 才恢复相同梯度。
若每个 micro-batch 单独裁剪，非线性的裁剪映射也不能穿过求和。
\end{proof}

生产实现采用 \path{loss_agg_mode=token-mean}，
在任何 loss 聚合之前填入当前全 batch 的有效 response token 数。
数据并行若采用梯度平均，每个 rank 的局部和还需乘以数据并行 rank 数，
使跨 rank 平均后恢复全局 token 均值。
每份 micro-batch 不能使用自身长度或上一 outer 的 token 数作为分母。
